\PassOptionsToPackage{heading=true}{ctex}
\documentclass[lang=cn,chinesefont=nofont,11pt,a4paper,bibend=bibtex]{elegantpaper}
\newcommand{\reportkind}{作品简介}
% Dududa submission reports. Upstream ElegantPaper is kept unmodified in vendor/.
\usepackage{tabularx,longtable,array,makecell}
\usepackage{fancyhdr,needspace,placeins,float}
\usepackage{tikz}
\usetikzlibrary{arrows.meta,positioning,fit,backgrounds,calc}
\usepackage{fvextra}
\usepackage{microtype}
\geometry{a4paper,left=25mm,right=25mm,top=25mm,bottom=24mm,headheight=15pt,headsep=8mm,footskip=11mm}
\setCJKmainfont{FandolSong-Regular.otf}[BoldFont=FandolSong-Bold.otf,ItalicFont=FandolKai-Regular.otf]
\setCJKsansfont{FandolHei-Regular.otf}[BoldFont=FandolHei-Bold.otf]
\setCJKmonofont{FandolFang-Regular.otf}
\setmonofont{lmmono10-regular.otf}[Scale=0.9]
\definecolor{ink}{HTML}{202D3A}
\definecolor{accent}{HTML}{284B69}
\definecolor{muted}{HTML}{596773}
\definecolor{rulegray}{HTML}{BCC6CE}
\definecolor{panel}{HTML}{F2F5F7}
\color{ink}
\hypersetup{unicode=true,colorlinks=true,linkcolor=accent,urlcolor=accent,citecolor=accent,pdfauthor={嘟嘟哒项目},pdfcreator={LaTeX / ElegantPaper}}
\linespread{1.28}
\setlength{\parskip}{0.25em}
\setlength{\emergencystretch}{2em}
\clubpenalty=10000
\widowpenalty=10000
\setlength{\tabcolsep}{5pt}
\renewcommand{\arraystretch}{1.32}
\newcolumntype{L}[1]{>{\raggedright\arraybackslash}p{#1}}
\newcolumntype{Y}{>{\raggedright\arraybackslash}X}
\BeforeBeginEnvironment{tabularx}{\par\noindent}
\AfterEndEnvironment{tabularx}{\par\addvspace{0.5\baselineskip}}
\ctexset{
  section={format={\Large\sffamily\bfseries\color{accent}},
    aftername={\hspace{0.65em}},beforeskip={2.0ex plus .5ex},afterskip=1.1ex},
  subsection={format={\large\sffamily\bfseries},
    aftername={\hspace{0.6em}},beforeskip={1.7ex plus .4ex},afterskip=0.7ex},
  subsubsection={format={\normalsize\sffamily\bfseries},aftername={\hspace{0.6em}}}
}
\setcounter{tocdepth}{2}
\setcounter{secnumdepth}{3}
\setlist[itemize]{leftmargin=2em,itemsep=0.25em,topsep=0.35em}
\setlist[enumerate]{leftmargin=2em,itemsep=0.25em,topsep=0.35em}
\captionsetup{font=small,labelfont=bf,labelsep=quad}
\pagestyle{fancy}
\fancyhf{}
\fancyhead[L]{\small\sffamily 嘟嘟哒}
\fancyhead[R]{\small\sffamily\reportkind}
\fancyfoot[L]{\footnotesize\color{muted}2026 年 9 月 · 技术材料}
\fancyfoot[R]{\small\thepage}
\renewcommand{\headrulewidth}{0.35pt}
\renewcommand{\footrulewidth}{0pt}
\fancypagestyle{plain}{\fancyhf{}\fancyfoot[R]{\small\thepage}\renewcommand{\headrulewidth}{0pt}}
\newcommand{\code}[1]{{\small\expandafter\nolinkurl\expandafter{\detokenize{#1}}}}
\newcommand{\source}[1]{\par\smallskip{\footnotesize\color{muted}依据：#1}\par}
\newcommand{\reportdate}{2026 年 9 月 6 日}
\newcommand{\repolink}{\href{https://github.com/ly382965/dududa}{github.com/ly382965/dududa}}
\DefineVerbatimEnvironment{Code}{Verbatim}{fontsize=\footnotesize,breaklines=true,breakanywhere=true,frame=single,rulecolor=\color{rulegray},framesep=3mm,xleftmargin=1mm,xrightmargin=1mm}
\tikzset{box/.style={draw=rulegray,fill=panel,rounded corners=1pt,align=center,minimum height=11mm,font=\small,inner sep=5pt},flow/.style={-{Stealth[length=2mm]},draw=accent,thick},future/.style={box,dashed,fill=white},labeltext/.style={font=\footnotesize,color=muted,align=center}}
\newcommand{\reportcover}[3]{%
  \begin{titlepage}
  \thispagestyle{empty}\color{ink}
  \noindent{\small\sffamily\color{accent}DUDUDA\quad /\quad TECHNICAL DOCUMENTATION}\par
  \vspace{22mm}
  \noindent{\fontsize{34}{44}\selectfont\sffamily\bfseries 嘟嘟哒\par}
  \vspace{4mm}
  \noindent{\fontsize{24}{32}\selectfont\sffamily #1\par}
  \vspace{8mm}\noindent{\color{accent}\rule{\textwidth}{1pt}}\par
  \vspace{6mm}\noindent{\Large 群聊原生社交智能体\par}
  \vspace{3mm}\noindent{\large 面向校园生活、教学支持与群体协作\par}
  \vspace{15mm}
  \noindent\begin{minipage}{0.94\textwidth}\normalsize #2\end{minipage}
  \vfill
  \noindent\begin{tabularx}{\textwidth}{@{}L{27mm}Y@{}}
  提交材料 & #3\\
  文档日期 & \reportdate\\
  项目入口 & \repolink\\
  实现依据 & 本地源码与已存验证记录；能力状态详见正文\\
  \end{tabularx}
  \vspace{12mm}
  \end{titlepage}
}

\linespread{1.2}
\renewcommand{\arraystretch}{1.15}
\setlength{\parskip}{0.15em}
\hypersetup{pdftitle={嘟嘟哒 作品简介},pdfsubject={作品背景、核心功能、模型调用与创新点}}
\begin{document}
\noindent{\small\sffamily\color{accent}DUDUDA\quad /\quad PROJECT OVERVIEW}\par
\vspace{5mm}
\noindent{\fontsize{27}{35}\selectfont\sffamily\bfseries 嘟嘟哒\quad 作品简介}\par
\vspace{3mm}
\noindent{\large 群聊原生社交智能体}\par
\vspace{2mm}
\noindent{\small 第 4 项提交材料\quad\textbar\quad \reportdate\quad\textbar\quad \repolink}\par
\vspace{3mm}\noindent{\color{accent}\rule{\textwidth}{0.7pt}}\par

\section{作品背景与解决的问题}
嘟嘟哒 面向 QQ 群中的持续交流，将自然语言理解、群聊参与、校园查询和人格表达组织成一套智能体系统。它以“生活在 QQ 群里的 AI 群友”为产品定位，首批服务于校园生活、教学支持和群体协作。

群聊中的请求经常分散在多人发言里：有人提出问题，有人补充条件，有人临时更正；课程、活动、通知和教务资料又来自不同网站。用户需要一个能够理解当前讨论、在适当时机参与、查询资料并说明来源的入口。管理员则需要按群控制参与频率和服务范围，查看回答所使用的模型、工具及运行结果。

作品围绕三个实际问题展开：\textbf{从连续讨论中理解需求；将需求落到可用的校园服务；让群聊参与有节奏、可配置、可解释。}明确提及 Bot 可发起任务；在目标群开启自动参与后，普通讨论也能形成回应机会。

\section{核心功能}
\begin{tabularx}{\textwidth}{@{}L{32mm}Y@{}}\toprule
功能 & 具体能力\\\midrule
群聊理解与参与 & 读取同群近期上下文，处理引用、追问与更正；按群配置参与概率、冷却和每小时上限。\\
评课与教学支持 & 查询课程、教师、评论与统计；区分同名条目，按讨论需求组织有来源的建议。\\
活动与教务 & 搜索第二课堂活动，读取详情和类别；查询学期、全校开课、考试与教学日历。\\
方案、通知与出行 & 查询不同年级的培养方案资料、校园通知及截止日期；按路线、日期和时段查校车。\\
自适应插件选择 & 在管理员已授权的只读插件集合中，根据多人连续讨论和当前请求选择、启用匹配的查询能力。\\
人格与扩展交互 & 独立 Persona 资产与表达规则；表情合成、复读和曲目信息等扩展由相应插件提供。\\
管理与调试 & Web 展示 Agent 对话、实际模型、能力与耗时，支持群配置、三档 Key 池和 MCP 连接检查。\\\bottomrule
\end{tabularx}

培养方案使用有版本的非官方研究快照，校车使用静态时刻表；二课认证在部署端配置。个人成绩、代报名、个人学时查询不属于本次校园智能体范围。

\clearpage
\section{技术架构与实现方法}
系统以 Python Runtime 为核心，通过 AstrBot/NapCat 接入 QQ，通过 MCP 与本地 Provider 连接校园资料，使用 Vue 3/Node.js 提供管理控制台。QQ 与 Web 预览复用同一运行核心。

\begin{figure}[H]\centering
\begin{tikzpicture}[node distance=6mm]
\node[box,text width=13cm] (a) {QQ / Web 入口\quad 消息身份、提及、引用与同群近期历史};
\node[box,text width=13cm,below=of a] (b) {上下文感知 $\rightarrow$ Social Engine $\rightarrow$ 模型路由};
\node[box,text width=13cm,below=of b] (c) {直接交流\quad 或\quad 能力选择 $\rightarrow$ 单步查询 $\rightarrow$ 结果校验};
\node[box,text width=13cm,below=of c] (d) {回答组织 $\rightarrow$ Persona 与格式检查 $\rightarrow$ QQ 回复 / Web 预览};
\draw[flow](a)--(b);\draw[flow](b)--(c);\draw[flow](c)--(d);
\end{tikzpicture}
\caption{当前主要运行流程。长期 Memory 与多步 Skill 另处于接入或后续建设阶段。}
\end{figure}

\textbf{规则与模型协同。}规则识别明确来源、平台信号和输入约束，模型提取意图与实体。结构化语义经过校验，再由确定性策略选择回应、查询、澄清或静默；身份、权限和输出目标由程序维护。

\textbf{业务能力统一。}课程、活动、通知等服务各自声明输入输出与 Provider，Planner 把自然语言条件绑定到有界业务查询。当前采用单步计划，MCP worker 使用独立解释器解决宿主与客户端的依赖版本差异。

\textbf{计算与表达分开。}模型档位、推理强度和正文篇幅分别设置。人格影响语气与信息呈现，不能改变工具事实、来源和用户指定的精确格式。

\clearpage
\section{WebUI 工作台设计}
WebUI 将真实 QQ 会话与智能体管理放在同一工作区：桌面端并列显示账号栏、会话列表、聊天窗和可折叠的 Agent 面板；手机端通过单页切换保留当前会话。界面提供浅色和深色主题，以明确的文字和状态色区分 QQ 在线、Runtime 可用及配置生效情况。

\begin{tabularx}{\textwidth}{@{}L{32mm}Y@{}}\toprule
页面区域 & 设计与使用方式\\\midrule
消息工作区 & 查看同群讨论、读取历史和搜索已缓存消息；账号与会话始终对应当前操作。\\
Agent 三个页签 & “对话”输入问题和阅读回答；“运行”检查实际模型、工具、上下文与耗时；“配置”维护群行为与插件范围。\\
模型与工具管理 & 管理三档 Key 池、MCP 连接和插件；分别显示保存结果及应用到 Runtime 的状态。\\
连接与异常反馈 & 首屏先显示工作区，后台刷新目录；错误指出历史读取、账号登录或 Runtime 调用等具体环节，并提供重试。\\\bottomrule
\end{tabularx}

浏览器以同源 HTTP/SSE 连接 Node.js 网关；网关负责 QQ 同步，并从服务端取得近期历史，交给 AstrBot 中的既有 Runtime 生成预览。预览可使用校园工具，结果留在 Web 中供管理员检查。该设计让管理员沿“阅读讨论—设置服务—检查回答”完成连续操作。配套《设计文档》设有独立的“WebUI 设计”一章，详述布局、数据流、交互及实现依据。

\section{创新点}
嘟嘟哒以\textbf{自适应}为核心：围绕当前群聊情境，分别调整使用的插件、计算路由和参与方式，并通过 WebUI 展示选择过程、支持人工配置。

\begin{enumerate}
\item \textbf{自适应插件选择。}结合讨论主题、多人参与、需求持续性与当前请求，在管理员授权的服务集合中选择并启用匹配插件。例如，选课讨论形成明确需求后启用评课查询，记录启用原因，将自然讨论连接到具体校园服务。
\item \textbf{任务驱动的自适应路由。}根据意图、复杂度信号、置信度、可用模型和预算选择计算档位，并独立配置推理强度与回答篇幅。简单交流与复杂分析可采用不同路由，实际选择随本轮任务变化，并保留选择依据。
\item \textbf{自适应群聊参与与表达。}结合提及、追问、持续讨论和群策略选择回复、查询、澄清或静默；再按任务组织回答篇幅与人格表达，使参与方式适应当前情境，同时遵守配置的冷却和频率限制。
\item \textbf{WebUI 工作台设计。}将真实 QQ 会话与 Agent 对话、运行、配置整合为连续操作；管理员按群设置偏好、允许范围与锁定项，观察实际模型、插件、上下文和耗时，形成“配置—预览—观察—调整”的自适应管理流程。
\end{enumerate}

这些创新集中在可解释、可配置的情境适应机制：当前以语义感知、规则与策略组合实现，在线学习另作为后续研究方向。

\clearpage
\section{使用的模型及 API 调用方式}
本次演示使用 DeepSeek，常规配置如下。界面和内部档位均为程序别名，实际模型以请求中的 model 字段为准。

\begin{table}[htbp]\centering
\caption{模型配置}
\begin{tabularx}{\textwidth}{@{}L{27mm}L{21mm}L{45mm}Y@{}}\toprule
界面档位 & 内部别名 & 实际模型 ID & 常规推理强度\\\midrule
Luna 快速 & haiku & \code{deepseek-v4-flash} & low\\
Terra 标准 & sonnet & \code{deepseek-v4-flash} & high\\
Sol 深度 & opus & \code{deepseek-v4-pro} & max\\\bottomrule
\end{tabularx}
\end{table}

服务地址为 \url{https://api.deepseek.com}，协议为 \textbf{OpenAI 兼容 Chat Completions}。调用链为 Runtime $\rightarrow$ AstrBotModelProviderAdapter $\rightarrow$ AstrBot Provider.text\_chat $\rightarrow$ DeepSeek HTTP API。后端传入提示词、模型、输出额度、thinking 与 reasoning\_effort，接收回答正文和 token 用量。

2026 年 9 月 6 日录制查询配置保留快速档轻量推理，标准档和深度档使用 \code{thinking: disabled}，直接组织工具资料。以下为标准档录制模式的等价请求示意：

\begin{Code}
POST https://api.deepseek.com/chat/completions
Authorization: Bearer <DEPLOYMENT_API_KEY>
Content-Type: application/json
{
  "model": "deepseek-v4-flash",
  "messages": [
    {"role": "system", "content": "<task instructions>"},
    {"role": "user", "content": "<question and tool evidence>"}
  ],
  "thinking": {"type": "disabled"},
  "max_tokens": 8192
}
\end{Code}

常规推理模式改为 thinking enabled，并设置 low/high/max 等 reasoning\_effort。真实 API Key 在部署端管理。校园工具由程序执行后将结果交给模型，当前主链不依赖模型原生 Function Calling 自主执行工具。

字段依据：\href{https://api-docs.deepseek.com/guides/thinking_mode/}{DeepSeek 官方 Thinking Mode 文档}，2026 年 9 月 6 日核对；项目实际参数以模型适配器及运行配置为准。

\section{实际使用示例}
\textbf{选课讨论：}群内连续提出讲解风格、证明训练和课业负担等需求；条件满足且已有授权后，系统启用评课能力，查询课程与教师评论，按不同需求组织建议并附链接。

\textbf{校园查询：}用户询问二课英语角、某学期开课记录、培养方案学分或校车时刻，系统定位业务、提取条件并返回资料。对不存在、未取得或版本不明确的信息，应说明范围或请求澄清。

\clearpage
\section{已有成果与验证}
当前成果包括 QQ 接入、近期上下文、Social Engine、按群自动参与、授权自适应启用、六类校园查询、三档模型配置、独立 Persona 资产和 Web 控制台。验证材料按执行层级分别统计：

\begin{tabularx}{\textwidth}{@{}L{32mm}Y@{}}\toprule
验证类型 & 已存结果与解释\\\midrule
100 条离线流程 & 94 条 Runtime 响应，6 条入口规则处理；187 次脚本模型调用、53 次 MCP fixture 调用、11 次本地能力调用。预期检查无不匹配，真实 QQ 发送为 0。本次报告编制已复跑。\\
26 项原子 MCP 查询 & 选定案例 26/26 通过；包括真实校园服务与可选服务的查询，证明这些调用可用。\\
5 项真实模型问答 & Flash 标准档预览 5/5 通过对应检查；涵盖开课、培养方案、通知、校车、二课，单次总耗时 10.92—29.07 秒。\\
四人讨论排练 & 20 条模拟发言，前 19 条评课关闭，第 20 条启用并查询；原生排练记录 1 次输出调用与完成处置。\\
广覆盖真实预览 & 较早 93 次真实模型预览审查为 60 通过、25 部分、8 失败；仍保留问题清单，未用后续少量样本代替全量复测。\\\bottomrule
\end{tabularx}

离线流程验证不等于真实模型准确率，选定演示样本也不构成总体性能或效果统计。模拟群成员输入与真实 QQ 用户完整发言分别说明；具体证据、复现命令和限制见配套《设计文档》的“验证结果与完成度”一章。

\section{当前边界与后续建设}
Memory v2 已实现作用域检索、中文 BM25、写入筛选和生命周期组件，主动调度、摘要与 Bandit 也有独立或离线成果。\textbf{长期记忆尚未接入正式聊天，通用多步 Skill、多 OC 切换、群友印象和在线学习仍需建设。}

近期优化聚焦查询前澄清、来源完整性、近期摘要的时间与发言归属，以及查询耗时和参与误触发。后续以实际群聊任务评估长期记忆和组合能力的收益，逐步扩展到更多校园及群体协作场景。

\section{交付与技术资料}
作品提供可部署源码、依赖锁文件、核心 wheel、Web 构建产物、Docker Compose 和测试入口。部署者配置模型服务、必要校园认证与 QQ 账号后可连接运行，也可先使用无需模型 Key 的离线流程。

配套设计文档整合本地技术资料，附来源索引和 LaTeX 源码。两份报告实际采用 \href{https://github.com/ElegantLaTeX/ElegantPaper}{ElegantPaper} 中文模板，原模板许可证随源码提供。

\noindent\textbf{项目入口：}\repolink\par
\noindent\textbf{材料依据：}主仓源码与各服务 README、录制验证记录、群聊 100 题结果、Memory/Persona 设计及离线评测数据卡；详细路径见配套设计文档与本地文档汇总。
\end{document}
